% Options for packages loaded elsewhere
\PassOptionsToPackage{unicode}{hyperref}
\PassOptionsToPackage{hyphens}{url}
\PassOptionsToPackage{dvipsnames,svgnames,x11names}{xcolor}
%
\documentclass[
]{article}
\usepackage{amsmath,amssymb}
\usepackage{iftex}
\ifPDFTeX
  \usepackage[T1]{fontenc}
  \usepackage[utf8]{inputenc}
  \usepackage{textcomp} % provide euro and other symbols
\else % if luatex or xetex
  \usepackage{unicode-math} % this also loads fontspec
  \defaultfontfeatures{Scale=MatchLowercase}
  \defaultfontfeatures[\rmfamily]{Ligatures=TeX,Scale=1}
\fi
\usepackage{lmodern}
\ifPDFTeX\else
  % xetex/luatex font selection
  \setmainfont[]{DejaVu Serif}
  \setsansfont[]{DejaVu Sans}
  \setmonofont[]{DejaVu Sans Mono}
\fi
% Use upquote if available, for straight quotes in verbatim environments
\IfFileExists{upquote.sty}{\usepackage{upquote}}{}
\IfFileExists{microtype.sty}{% use microtype if available
  \usepackage[]{microtype}
  \UseMicrotypeSet[protrusion]{basicmath} % disable protrusion for tt fonts
}{}
\makeatletter
\@ifundefined{KOMAClassName}{% if non-KOMA class
  \IfFileExists{parskip.sty}{%
    \usepackage{parskip}
  }{% else
    \setlength{\parindent}{0pt}
    \setlength{\parskip}{6pt plus 2pt minus 1pt}}
}{% if KOMA class
  \KOMAoptions{parskip=half}}
\makeatother
\usepackage{xcolor}
\usepackage[margin=0.7in]{geometry}
\setlength{\emergencystretch}{3em} % prevent overfull lines
\providecommand{\tightlist}{%
  \setlength{\itemsep}{0pt}\setlength{\parskip}{0pt}}
\setcounter{secnumdepth}{-\maxdimen} % remove section numbering
\ifLuaTeX
\usepackage[bidi=basic]{babel}
\else
\usepackage[bidi=default]{babel}
\fi
\babelprovide[main,import]{russian}
\ifPDFTeX
\else
\babelfont{rm}[]{DejaVu Serif}
\fi
% get rid of language-specific shorthands (see #6817):
\let\LanguageShortHands\languageshorthands
\def\languageshorthands#1{}
\ifLuaTeX
  \usepackage{selnolig}  % disable illegal ligatures
\fi
\usepackage{bookmark}
\IfFileExists{xurl.sty}{\usepackage{xurl}}{} % add URL line breaks if available
\urlstyle{same}
\hypersetup{
  pdflang={ru-RU},
  colorlinks=true,
  linkcolor={blue},
  filecolor={Maroon},
  citecolor={Blue},
  urlcolor={blue},
  pdfcreator={LaTeX via pandoc}}

\author{}
\date{}

\begin{document}

\section{1. Идея и цель
проекта}\label{ux438ux434ux435ux44f-ux438-ux446ux435ux43bux44c-ux43fux440ux43eux435ux43aux442ux430}

\subsection{Идея
проекта}\label{ux438ux434ux435ux44f-ux43fux440ux43eux435ux43aux442ux430}

Проект представляет собой программную игру-симуляцию, в которой
несколько игроков управляют инвестиционным капиталом и принимают решения
о вложении денег в строительство жилых домов и супермаркетов.

Игроки развивают объекты в одном микрорайоне и конкурируют между собой
за получение максимальной прибыли. На результат игры влияют стоимость и
сроки строительства, спрос на жильё, продажи в супермаркетах, цены на
жильё и затраты на рекламу.

Каждый шаг игры соответствует одному месяцу. В течение игры игроки
должны распределять имеющийся капитал между различными объектами и
принимать решения, влияющие на их дальнейший доход.

\subsection{Цель
проекта}\label{ux446ux435ux43bux44c-ux43fux440ux43eux435ux43aux442ux430}

Цель проекта --- разработать игровую программную систему, моделирующую
процесс инвестирования в строительство и позволяющую игрокам принимать
экономические решения в условиях изменяющегося спроса и ограниченного
капитала.

Победителем становится игрок, который к концу установленного количества
игровых месяцев имеет наибольший общий капитал.

\subsection{Основная идея игрового
процесса}\label{ux43eux441ux43dux43eux432ux43dux430ux44f-ux438ux434ux435ux44f-ux438ux433ux440ux43eux432ux43eux433ux43e-ux43fux440ux43eux446ux435ux441ux441ux430}

В начале игры каждый игрок получает одинаковый начальный капитал. Затем
игроки выбирают, в какие объекты и в каком количестве инвестировать
средства.

В игре предусмотрены два основных типа объектов:

\begin{itemize}
\tightlist
\item
  жилые дома;
\item
  супермаркеты.
\end{itemize}

Строительство объектов требует определённых затрат и занимает несколько
месяцев. При этом строительство объектов одного типа может влиять на
эффективность объектов другого типа: дополнительные супермаркеты могут
повышать спрос на расположенное рядом жильё, а дополнительные дома могут
увеличивать уровень продаж в соседних супермаркетах.

Кроме строительства, игроки принимают решения о продаже жилья, его цене
и затратах на рекламу.

\subsection{Основные участники
системы}\label{ux43eux441ux43dux43eux432ux43dux44bux435-ux443ux447ux430ux441ux442ux43dux438ux43aux438-ux441ux438ux441ux442ux435ux43cux44b}

В проекте выделяются следующие основные элементы:

\begin{itemize}
\tightlist
\item
  игроки
\item
  строительные объекты
\item
  рынок жилья
\item
  риэлторское агентство
\item
  супермаркеты
\item
  игровая система, управляющая ходом моделирования
\end{itemize}

Риэлторское агентство обеспечивает продажу выставленного игроками жилья
и распределяет удовлетворяемый спрос в зависимости от заявленной цены,
рекламы и других факторов.

\subsection{Конечный
результат}\label{ux43aux43eux43dux435ux447ux43dux44bux439-ux440ux435ux437ux443ux43bux44cux442ux430ux442}

Игра продолжается заданное количество месяцев. После завершения игры
рассчитывается общий капитал каждого игрока, включающий стоимость его
объектов и выручку от продаж.

Победителем считается игрок, обладающий наибольшим общим капиталом.

\end{document}
